\chapter{Convergencia débil}

\begin{definicion}
    Sea $X$ un espacio normado y $\{x_n\} \subseteq X$. Se dice que $x_n$ \textbf{converge débilmente} a $x \in X$ si
    \begin{equation*}
        F(x_n) \longrightarrow F(x) \quad \forall F \in X^*
    \end{equation*}
    y escribiremos $x_n \rightharpoonup x$. \\
\end{definicion}

\begin{observacion}
    Llamaremos \tbf{convergencia fuerte} a la convergencia habitual, es decir, $x_n \longrightarrow x$ si \mbox{$\|x_n - x\| \longrightarrow 0$}.
\end{observacion}

\begin{prop}
    El límite débil es único.
\begin{proof}
    Sean $x, y \in X$ tales que $x_n \rightharpoonup x$ y $x_n \rightharpoonup y$. Entonces \mbox{$F(x_n) \longrightarrow F(x) = F(y) \quad \forall F \in X^*$}, de lo cual:
    \begin{equation*}
        F(x-y) = 0 \quad \forall F \in X^* \implies \|x-y\| \overset{(*)}{=} \sup_{\|F\|=1, F\in X^*} |F(x-y)| = 0 \implies x = y.
    \end{equation*}
    donde en $(*)$ hemos usado el Corolario \ref{coro:hahn-banach-analitico-norma} del teorema de Hahn-Banach. \\
\end{proof}
\end{prop}

\begin{prop}
    Si $x_n \longrightarrow x$ en $X$ (convergencia fuerte), entonces $x_n\rightharpoonup x$ (convergencia débil).

\begin{proof}
    Sea $F\in X^*$, entonces:
    \begin{equation*}
        \|F(x_n) - F(x)\| \le \|F\| \cdot \underbrace{\|x_n - x\|}_{\to 0} \longrightarrow 0
    \end{equation*}
    De lo cual $F(x_n) \longrightarrow F(x) \quad \forall F \in X^*$.
\end{proof}
\end{prop}

\observacion{
    El recíproco es falso en general, es decir, $x_n \rightharpoonup x$ no implica $x_n \longrightarrow x$. De hecho, si $X = l_2$ y $e_n$ es la sucesión canónica de $l_2$, entonces $e_n \rightharpoonup 0$ pero $\|e_n\| = 1 \not\to 0$.
}

\begin{prop}
    % //TODO: entender demostración
    Sea $\{x_n\}$ una sucesión convergente débilmente a $x \in X$. Entonces $\{x_n\}$ es acotada y se cumple
    \begin{equation*}
        \|x\| \le \liminf_{n \to +\infty} \|x_n\|
    \end{equation*}
\begin{proof}
    Sea $J: X \longrightarrow X^{**}$ el mapa canónico ($x \mapsto J_x$), entonces si $x_n \rightharpoonup x$ tenemos que $F(x_n) \to F(x) \quad \forall F \in X^*$, es decir,
    $J_{x_n} \longrightarrow J_x$ en $X^{**}$ es puntualmente convergente $\xrightarrow{PUL}$ ($J_{x_n}$) es uniformemente acotada: $\exists M > 0 :$
    \begin{equation*}
        \|J_{x_n}\| \le M \quad \forall n \implies \|x_n\| \le M \quad \forall n \implies \{x_n\} \text{ es acotada, de lo cual:}
    \end{equation*}

    $\forall F \in X^* \quad |F(x_n)| \le \|F\| \cdot \|x_n\|$. Tomamos $\|F\| \le 1$ entonces $|F(x_n)| \le \|x_n\|$ para todo $n$, por lo que $|F(x)| = |\lim_{n} F(x_n)| = |\liminf_{n} F(x_n)| \le \liminf_n \|x_n\|$ para todo $F$ con $\|F\| \le 1$.
    \begin{equation*}
        \sup_{\|F\| \le 1} |F(x)| \le \liminf_{n \to +\infty} \|x_n\| \implies \|x\| \le \liminf_{n \to +\infty} \|x_n\| \quad \square
    \end{equation*}
\end{proof}
\end{prop}

\begin{definicion}
    Sea $C\subseteq X$ se dice \tbf{débilmente cerrado} (por sucesiones) si para toda $\{x_n\} \subseteq C$ convergente débilmente a $x \in X$, entonces $x \in C$.
\end{definicion}

\begin{definicion}
    Sea $K\subseteq X$ se dice \tbf{débilmente compacto} (por sucesiones) o \tbf{débilmente secuancialmente compacto} si para toda $\{x_n\} \subseteq K$ existe una subsucesión convergente débilmente en $K$.
\end{definicion}

\begin{observacion}
    La \tbf{envolvente convexa} (o convex hull) de un conjunto de puntos $\{x_n\}$ es el conjunto más pequeño que contiene a todos esos puntos y que además es convexo.
\end{observacion}

\begin{teo}[Mazur]
    Sea $\{x_n\} \subseteq X$ tal que $x_n \rightharpoonup x_0$, entonces $x_0 \in \overline{Co(x_n)}$, donde:
    \begin{equation*}
        Co(x_n) := \text{envolvente convexa de } \{x_n\}
    \end{equation*}
\begin{proof}
    Supongamos por reducción al absurdo (P.A.) que $x_0 \notin \overline{Co(x_n)}$ : $d(x_0, Co) = \delta > 0$, entonces $B_\delta(x_0) \cap \overline{Co(x_n)} = \emptyset$ y, por lo tanto, por el teorema de Hahn-Banach geométrico:
    \begin{equation*}
        \exists F \in X^*, \exists \alpha \in \mathbb{R} \text{ tal que } F(x_0) < \alpha < F(x_n) \quad \forall n
    \end{equation*}
    pero $F(x_n) \longrightarrow F(x_0)$, lo cual es un ABSURDO.
\end{proof}
\end{teo}

\begin{coro}
    Sea $C \subseteq X$ cerrado y convexo; entonces es débilmente cerrado.
\begin{proof}
    Por el teorema anterior, dada $\{x_m\} \subseteq C$ tal que $x_m \rightharpoonup x$, entonces $x \in \overline{Co(x_m)} \subseteq C \implies$ es débilmente cerrado. 
\end{proof}
\end{coro}

\observacion{
    Sea $X$ de dimensión finita, entonces la convergencia débil implica la convergencia fuerte. 
}

\begin{definicion}
    Sean $X, Y$ espacios normados de Banach, $T$ se dice \tbf{débilmente secuancialmente continuo} si para cada $\{x_n\} : x_n \rightharpoonup x$, entonces $T(x_n) \rightharpoonup T(x)$.
\end{definicion}

\begin{prop}
    Sea $T: X \longrightarrow Y$ lineal, entonces
    \begin{equation*}
        T \text{ es débilmente secuancialmente continuo} \iff T \text{ es continuo}
    \end{equation*}
\begin{proof}
    \begin{description}
        \item[$\Longrightarrow)$] $T$ es débilmente continuo, entonces si $x_m \rightharpoonup x$ se tiene que $T(x_m) \rightharpoonup y$ con $T(x) = y$ por la unicidad del límite. \\
        \noindent
        Si $T$ es débilmente continuo, entonces $T$ es cerrada; de hecho:
        Si $x_m \to x$ y $T(x_m) \to y$, entonces $\begin{cases} x_m \rightharpoonup x \\ T(x_m) \rightharpoonup y \end{cases} \implies T(x) = y$. \\
        \noindent
        Por el teorema del \textbf{gráfico cerrado}, $T$ es continua.
        \item[$\Longleftarrow)$] Sea $\{x_m\} \in X : x_m \rightharpoonup x$, sea $G \in Y^*$, entonces: $X \xrightarrow{T} Y \xrightarrow{G} \mathbb{R}$ con $T, G$ continuas $\implies G \circ T$ continua.
        \[ \therefore G \circ T \in X^* \quad \text{y por lo tanto} \quad (G \circ T)(x_m) \longrightarrow (G \circ T)(x) \quad \forall G \in Y^* \implies T(x_m) \rightharpoonup T(x) \]
    \end{description}
\end{proof}
\end{prop}

\begin{ejercicio}
    Sea $\{F_n\} \subseteq X^*$ tal que $F_n \to F$ en $X^*$ y supongamos que existe $\{x_n\} : x_n \rightharpoonup x$, entonces $F_n(x_n) \longrightarrow F(x)$. \\
    \noindent
    Véamoslo
    \begin{equation*}
        \|F_n(x_n) - F(x)\| \le \|F_n(x_n) - F_n(x)\| + \|F_n(x) - F(x)\| \le \|F_n - F\| \cdot \underbrace{\|x_n\|}_{\text{acotada}} + \underbrace{|F_n(x) - F(x)|}_{\to 0} \longrightarrow 0
    \end{equation*}
\end{ejercicio}

\begin{teo}[Eberlein]
    Si $X$ es un espacio reflexivo, entonces toda sucesión acotada admite una subsucesión débilmente convergente. \textit{\small(Los conjuntos acotados son débilmente secuencialmente compactos)}.

\begin{proof}
    % //TODO: leer demostración
    Sea $\{x_m\} \subseteq X$ y sea $Y = \overline{\text{span}(\{x_m\})}$ el subespacio de $X$ separable por construcción y reflexivo por ser cerrado. Como $Y$ es reflexivo, tenemos que la inyección canónica
    \Func{J}{Y}{Y^{**}}{y}{J_y}
    es un isomorfismo, por lo que $Y^{**}$ es reflexivo y aplicando el Teorema \ref{teo:separable-dual-separable} se tiene que $Y^*$ es separable, por lo que existe una sucesión $\{F_n\}$ densa en $Y^*$. \\
    \noindent
    Sea una sucesión $\{x_n\}$ acotada en $Y$, es decir, $\exists M > 0 : \|x_n\| \le M$ para todo $n$. Como $\{F_1(x_n)\}_n$ es acotada en $\mathbb{R}$, ya que $F_1\in Y^*$ y por tanto es lineal y acotada, entonces
    \begin{equation*}
        |F_1(x_n)| \le \|F_1\| \cdot \|x_n\| \le M \cdot \|F_1\| \quad \forall n
    \end{equation*}
    y por tanto admite una subsucesión convergente en $\mathbb{R}$ que es $\{F_1(x_n^1)_n\}$. Iterando obtenemos que 
    \begin{equation*}
        \{F_2(x_n^1)\}_n \text{ es acotada en } \mathbb{R} \implies \exists \{F_2(x_n^2)\}_n \text{ subsucesión convergente en } \mathbb{R}
    \end{equation*}
    y si seguimos iterando obtenemos
    \begin{equation*}
        \{x_m^1\}, \{x_m^2\}, \dots, \{x_m^k\}, \dots \text{ sucesiones  en } X \text{ tales que } \{F_i(x_m^j)\} \text{ es convergente en } \mathbb{R} \text{ para } i \le j
    \end{equation*}
    donde lo que hemos utilizado por ejemplo para $i=1$ y $j=2$, es que $\{F_1(x_n^1)\}_n$ es convergente y $\{F_1(x_n^2)\}_n$ es una subsucesión de $\{F_1(x_n^1)\}_n$ por ser $(x_n^2)$ subsucesión de $(x_n^1)_n$, por lo que $\{F_1(x_n^2)\}_n$ es convergente al mismo límite. \\
    
    \noindent
    De esta forma tenemos que en general, la subsucesión del nivel $j$, llamada $\{x_n^j\}$, tiene la propiedad de que hace converger a todos los funcionales anteriores $\{F_1, F_2, \dots, F_j\}$. Veámos ahora que $\{F_i(x_k^k)\}_k$ es convergente en $\mathbb{R} \quad \forall i \ge 1$, analicemos los términos de la diagonal a partir del índice $k = i$, puesto que para los términos anteriores a $i$ no se tiene la propiedad de convergencia, pero da igual pues no afectan al límite.
    \begin{itemize}
        \item El término $x_i^i$ pertenece a la subsucesión del nivel $i$.
        \item El término $x_{i+1}^{i+1}$ pertenece a la subsucesión del nivel $i+1$, la cual es una subsucesión de la del nivel $i$.
        \item En general para todo $k \ge i$, el término $x_k^k$ es un elemento que pertenece a la subsucesión del nivel $i$, por lo que $\{F_i(x_n^k)\}_n$ es convergente, de hecho, $\{F_i(x_k^k)\}_k$ es convergente. \\
    \end{itemize}
    Como $\{F_i\}_i$ es denso en $Y^*$, entonces $\{F(x_k^k)\}_k$ es convergente $\forall F \in Y^*$, veámoslo

    \[ |F(x_k^k) - F(x_j^j)| \le |F(x_k^k) - F_i(x_k^k)| + |F_i(x_k^k) - F_i(x_j^j)| + |F_i(x_j^j) - F(x_j^j)| \le \]
    \[ \le \underbrace{\|F - F_i\|}_{\le \varepsilon} \cdot \underbrace{\|x_k^k\|}_{\le M} + |F_i(x_k^k) - F_i(x_j^j)| + \underbrace{\|F - F_i\|}_{\le \varepsilon} \cdot \underbrace{\|x_j^j\|}_{\le M} \le 2\varepsilon M + |F_i(x_k^k) - F_i(x_j^j)| \]
    Fijado $\varepsilon > 0$, sea $\|F - F_i\| < \varepsilon$ y $\exists M > 0 : \|x_k^k\| \le M$, entonces como $(F_i(x_k^k))$ es convergente, tenemos que $(F_i(x_k^k))$ es de Cauchy, por lo que $|F_i(x_k^k) - F_i(x_j^j)| < \varepsilon$, de lo cual
    \begin{equation*}
        |F(x_k^k) - F(x_j^j)| \le 2\varepsilon M + \varepsilon \xrightarrow{\varepsilon \to 0} 0
    \end{equation*}
    Por lo tanto $\{F(x_k^k)\}_k$ es convergente; demostremos que converge a $F(x_0)$. \\

    \noindent
    Sea $J: Y \to Y^{**}$ y $(J_{x_k^k})_k \in Y^{**} : J_{x_k^k}(F) = F(x_k^k) \quad \forall F \in Y^*$. $J_{x_k^k}(F)$ es puntualmente convergente a un funcional  $G(F) = \lim\limits_{k \to \infty} J_{x_k^k}(F)$. Por un corolario del teorema de Banach-Steinhaus, $G \in Y^{**}$ (funcional lineal y acotado) y puesto que $X$ es reflexivo, $\exists x_0 \in Y : G = J_{x_0}$, entonces
    \[ J_{x_k^k}(F) = F(x_k^k) \xrightarrow{k \to +\infty} G(F)=J_{x_0}(F)=F(x_0)  \implies x_k^k \rightharpoonup x_0 \]
\end{proof}
\end{teo}

\begin{coro}
    Sea $C$ cerrado, convexo y acotado en $X$, un espacio de Banach reflexivo. Entonces $C$ es débilmente secuencialmente compacto.
\begin{proof}
    Sea $C$ acotado, entonces si $\{x_m\} \subseteq C$, entonces $\exists x_0 \in X : x_m \rightharpoonup x_0 \in \overline{Co(C)} = C$ y por lo tanto se sigue el resultado.
\end{proof}
\end{coro}

\observacion{
    También vale el recíproco (difícil), es decir, si $C$ es débilmente secuancialmente compacto, entonces $C$ es acotado.
}

\begin{coro}
    Sea $X$ un espacio de Banach reflexivo y sea $C$ cerrado no vacío y convexo; entonces $C$ admite un punto de norma mínima.
\begin{proof}
    Sea $d := \inf_{x \in C} \|x\|$ y sea $\{x_m\} \in X$ tal que $\|x_m\| \longrightarrow d$, entonces $\{x_m\}$ es acotada y, por lo tanto, $\exists \{x_{m_k}\} :$
    \[ x_{m_k} \rightharpoonup x_0 \in C \quad (\text{por ser } C \text{ cerrado y convexo}). \quad \text{De lo cual:} \]

    $\|x_0\| \le \liminf_k \|x_{m_k}\| = d$, pero $d = \inf_{x \in C} \|x\|$, por lo que $\|x_0\| = d$.
\end{proof}
\end{coro}

\begin{coro}
    Sea $C$ cerrado y convexo no vacío, entonces 
    $$\forall x \in X \quad \exists y \in C : \|x - y\| = \min_{z \in C} \|x - z\| = d(x, C)$$
\end{coro}

\begin{ejemplo}
    $x_m \rightharpoonup x$ en $\ell_p$ con $1<p<\infty \iff \{x_m\}$ es acotada y $x_m^j \longrightarrow x^j \quad \forall j \ge 1$ \\
    \begin{description}
        \item[$\Longrightarrow)$] Sea $F \in \ell_p^* \cong \ell_q$ con $\frac{1}{p} + \frac{1}{q} = 1$, entonces:
        $$F(y) = y^j \implies F(x_m) = x_m^j \longrightarrow x^j = F(x)$$
        \item[$\Longleftarrow)$] $\dots$
    \end{description}
\end{ejemplo}

\observacion{
    En $\ell_1$ no vale, pues tomando $\{e_n\} \subseteq \ell_1$, $e_n=(0, 0, \dots, 0, \underbrace{1}_{n}, 0, \dots)$, sabemos que $(e_n)$ está acotada, entonces $e_n^j \longrightarrow 0$ para todo $j$, pero $e_n \not\rightharpoonup 0$. De hecho, \\

    \noindent
    $\ell_1^* \cong \ell_\infty$ y sea $y = (1, 1, \dots, 1) \in \ell_\infty$: $F_y(e_m) = \sum_j e_m^j y_j = \sum_{j=1}^\infty e_m^j = 1 \not\to 0$.
}

\begin{ejemplo}
    $f_m \rightharpoonup f$ en $L^p \iff \int f_m g \longrightarrow \int f g \quad \forall g \in L^q$ por el teorema de Riesz.
\end{ejemplo}

\begin{ejemplo}
    $C^0([0,1])$ : sea $\{x_m\} \in C^0([0,1])$, entonces: 
    $x_m \rightharpoonup x \iff \begin{cases} \{x_m\} \text{ es acotada en } \|\cdot\|_\infty \\ x_m(t) \longrightarrow x(t) \quad \forall t \in [0,1] \end{cases}$

    $\Rightarrow$ Sea $F_t \in C^0([0,1])^*$ tal que $F_t(x) = x(t)$, entonces: 
    $x_m \rightharpoonup x \Rightarrow F_t(x_m) \longrightarrow F_t(x) \iff x_m(t) \longrightarrow x(t)$

    \vspace{1em}

    $\Leftarrow$ Sea $F \in C^0([0,1])^*$, entonces $F$ se representa con funciones de variación acotada $g$:

    $F(x) = \int_0^1 x(t) \, dg$ con $\{x_m\}$ acotada en $\|\cdot\|_\infty$, por lo que $\exists M > 0 : \|x_m\| \le M$ y $x_m(t) \longrightarrow x(t)$.

    $F_g(x_m) = \int_0^1 x_m(t) \, dg \xrightarrow{TCD} \int_0^1 x(t) \, dg = F_g(x) \Rightarrow x_m \rightharpoonup x$

\end{ejemplo}

\section{Convergencia $*-$débil}

\noindent
Sea $X$ un espacio normado y sean $\{f_n\} \in X^*$, estudiamos la convergencia débil en $X^*$

\begin{equation*}
    f_n \rightharpoonup  f \text{ en } X^* \text{ si } G(f_n) \longrightarrow G(f) \quad \forall G \in X^{**}
\end{equation*}
donde $J: X \longrightarrow X^{**}$, $J_x(f) = f(x)$ no es sobreyectiva en general y $J(X) \subseteq X^{**}$.

\begin{definicion}
    Sea $\{f_n\} \subseteq X^*$, diremos que $f_n$ \textbf{converge *-débilmente} a $f \in X^*$ si
    \begin{equation*}
        J_x(f_n) \longrightarrow J_x(f) \quad \forall x \in X
    \end{equation*}
    es decir, $f_n(x) \longrightarrow f(x) \quad \forall x \in X$. Se nota como 
    \begin{equation*}
        f_n \overset{*}{\rightharpoonup} f \text{ en } X^*
    \end{equation*}
\end{definicion}

\observacion{
    El límite $*-$débil de $\{f_n\}$ es único, pues si $f_n \overset{*}{\rightharpoonup} f_1$ y $f_n \overset{*}{\rightharpoonup} f_2$, entonces \mbox{$f_n(x) \longrightarrow f_1(x)$} y \mbox{$f_n(x) \longrightarrow f_2(x)$} puntualmente para todo $x \in X$, por lo que $f_1(x) = f_2(x)$ para todo $x \in X$.
}

\observacion{
    Si $f_n \longrightarrow f$ en $X^*$ (convergencia fuerte), entonces $f_n \rightharpoonup f$ en $X^*$ (convergencia débil), entonces $f_n \overset{*}{\rightharpoonup} f$ en $X^*$.
}
\begin{prop}
    La convergencia fuerte implica la convergencia débil, y esta última implica la convergencia $*-$débil.
\end{prop}

\begin{ejemplo}
    $l_0 \subseteq \ell_\infty$ con $c_0 = \{ (x_m) \in \ell_\infty \mid x_m \to 0 \}$ y sea $F_n(x) = x_n$, entonces:

    $F_n \xrightarrow{w^*} 0$ porque $F_n(x) = x_n \longrightarrow 0$.

    $l_0^* \cong \ell_1 \implies \{F_n\}$ se representa con $e_n = (0, 0, \dots, 1, 0, \dots) \in \ell_1$. Veamos que $F_n \not\rightharpoonup 0$ en $l_0^*$:

    sea $G_y \in l_0^{**} = \ell_\infty$ con $y = (1, 1, 1, \dots, 1, \dots) \in \ell_\infty$, entonces:
    \[ G_y(F_n) = F_n(y) = 1 \xrightarrow{\forall n, n \to +\infty} 1 \neq 0 \]
    Por lo tanto, si $X$ no es reflexivo, la convergencia *-débil y la débil en $X^*$ son diferentes.
\end{ejemplo}

\observacion{
    Si $X$ reflexivo, entonces
    \begin{equation*}
        f_n \rightharpoonup f \text{ en } X^* \iff f_n \overset{*}{\rightharpoonup} f \text{ en } X^*
    \end{equation*}
    Si $X$ no es reflexivo, entonces la convergencia débil y la convergencia $*-$débil en $X^*$ son diferentes.
}

\begin{ejercicio}
    Sea $X$ un espacio de Banach y $\{F_n\} \subseteq X^*$ demostrar que
    \begin{itemize}
        \item [i)] $F_n \overset{*}{\rightharpoonup}F \implies (\|F_n\|)_n$ es acotada y $\|F\| \le \liminf_n \|F_n\|$
        \item [ii)] $F_n \overset{*}{\rightharpoonup} F$, $x_n \to x \implies F_n(x_n) \to F(x)$
    \end{itemize} 
    \begin{description}
        \item[i).] $F_m \xrightarrow{w^*} F \implies J_{F_m}(x) \xrightarrow{m \to +\infty} J_F(x) \quad \forall x \in X$, por lo tanto $F_m(x) \to F(x) \quad \forall x \in X$. Por el \textbf{PUL} (Principio de Uniforme Limitación), $\{F_m\}$ es acotada $\implies \|F_m\| \le M$.
    
        Además: $|F_m(x)| \le \|F_m\| \cdot \|x\| \quad \forall m, \forall x \in X \implies$ aplicamos el $\liminf$:
        
        $\liminf_m |F_m(x)| \le \liminf_m \|F_m\| \cdot \|x\| \quad \forall x \in X \implies$ aplicamos el $\sup_{\|x\|=1}$ con $\liminf_m |F_m(x)| = |F(x)|$:
        
        $\sup_{\|x\|=1} |F(x)| \le \liminf_m \|F_m\| \implies \|F\| \le \liminf_{m \to +\infty} \|F_m\|$

        \item[ii).] $|F_m(x_m) - F(x)| \le |F_m(x_m) - F_m(x)| + |F_m(x) - F(x)| \le \|F_m\| \cdot \|x_m - x\| \le M \cdot \|x_m - x\| \longrightarrow 0$
    
        (Donde se usa que $F_m \xrightarrow{w^*} F$ para el segundo término y que $\{F_m\}$ es acotada por $M$).
        
    \end{description}

\end{ejercicio}

\observacion{
    $F_m \overset{*}{\rightharpoonup} F$ no implica $F\in \overline{Co(F_m)}$.
}

\begin{ejercicio}
    Sea $X$ un espacio de Banach y sea $Y$ denso en $X$. Sea $\{F_n\} \subseteq X^*$, entonces
    \begin{equation*}
        F_n \overset{*}{\rightharpoonup} F \iff \begin{cases} \{F_n\} \text{ es acotada en } X^* \\ \{F_n(y)\} \text{ es convergente } \forall y \in Y \end{cases}
    \end{equation*}
    \begin{description}
        \item[$\Longrightarrow)$] Ya lo sabemos (por el Principio de Uniforme Limitación).
        \item[$\Longleftarrow)$] $F_n \overset{*}{\rightharpoonup} F$ si $F_n(x) \longrightarrow F(x) \quad \forall x \in X$. Sea $x \in X$ fijado, entonces \mbox{$\exists y \in Y : \|y - x\|_X < \varepsilon \quad \forall \varepsilon > 0$}.
        \[ |F_n(x) - F(x)| \le |F_n(x) - F_n(y)| + \underbrace{|F_n(y) - F(y)|}_{\to 0} + |F(y) - F(x)| \]
        como $\|F_n\| \le M$, entonces:
        \[ \le M \|x - y\| + \|F\| \cdot \|y - x\| \le (M + \|F\|) \varepsilon \longrightarrow 0 \]
    \end{description}
\end{ejercicio}

\begin{teo}
    Sea $X$ un espacio de Banach separable. Toda sucesión de funcionales $\{f_n\} \subseteq X^*$ acotada en $X^*$ admite una subsucesión $*$-débilmente convergente.
\begin{proof}
    % //TODO: leer demostración
    Sea $X = \{x_n\}$ una sucesión densa en $X$ y sea $\{f_m\} \subseteq X^* : \|f_m\| \le M \quad \forall m$. Entonces:
    $\{f_m(x_1)\}$ es acotada en $\mathbb{R}$, de hecho: $\|f_m(x_1)\| \le \|f_m\| \cdot \|x_1\| \le M \cdot \|x_1\|$, de lo cual:

    $\exists \{f_m^{(1)}(x_1)\} \subseteq \mathbb{R}$ convergente. Análogamente, $\{f_m^{(1)}(x_2)\} \subseteq \mathbb{R}$ es acotada:
    $\|f_m^{(1)}(x_2)\| \le M \cdot \|x_2\| \implies \exists \{f_m^{(2)}(x_2)\}$ convergente en $\mathbb{R}$. Iterando:

    existe $\{f_n^{(k)}(x_k)\}$ sucesión convergente subsucesión de $(f_n^{(k-1)}(x_k))$.

    Sea $\{f_k^{(k)}\}_k$ la sucesión que tiene la propiedad: $\{f_k^{(k)}(x_j)\}$ converge para cada $j \ge 1$.
    Puesto que $\{x_j\}$ es denso en $X$, se sigue que $\{f_k^{(k)}(x)\}$ converge $\forall x \in X$, de hecho: $\forall j \ge 1$

    \vspace{1em}
    $|f_k^k(x) - f_m^m(x)| \le |f_k^k(x) - f_k^k(x_j)| + |f_k^k(x_j) - f_m^m(x_j)| + |f_m^m(x_j) - f_m^m(x)| \le$
    \[ \le \|f_k^k\| \cdot \|x - x_j\| + |f_k^k(x_j) - f_m^m(x_j)| + \|f_m^m\| \cdot \|x - x_j\| \le 2M\varepsilon + \varepsilon \longrightarrow 0 \]

    $\implies \{f_k^k(x)\}$ es de Cauchy en $\mathbb{R} \implies$ es convergente a $f(x) : f_k^k(x) \longrightarrow f(x)$ con $f(x)$ lineal y acotado por PUL. $\quad \square$

\end{proof}
\end{teo}

\observacion{
    Si $X$ no es seperable, en general no vale.
}

\begin{ejemplo}
    Sea $X=l_\infty$ no es separable, veámos que no se cumple entonces. \\
    \noindent
    $\{f_m\} \in \ell_\infty^* \quad f_m(x) = x_m$ con $x \in \ell_\infty \implies \|f_m\| = 1$ es acotada pero no tiene subsucesiones que converjan débilmente. De hecho, si $\{f_{m_j}\}$ es una subsucesión, tomamos 
    $$\tilde{x} = (\tilde{x}_k), \ \tilde{x}_k = \begin{cases} 1 & k = m_j \text{ par} \\ 0 & k = m_j \text{ impar} \end{cases}$$
    entonces $f_{m_j}(\tilde{x}) = \tilde{x}_{m_j}$ no es convergente.
\end{ejemplo}

\begin{teo}[Teorema de Banach-Alaoglu]
    Sea $X$ un espacio de Banach, entonces $B^*$, la bola cerrada en $X^*$, es débilmente-* compacta.
\end{teo}

\observacion{
    En un espacio de Hilbert (lo veremos más adelante), la convergencia débil y la convergencia débil$-*$ son exactamente la misma cosa, pues $\mathcal{H}^* \cong \mathcal{H}$.
}

\begin{coro}
    Si $H$ es un espacio de Hilbert, la bola cerrada $B = \{x \in H : \|x\| \le 1\}$ es débilmente compacta. Por tanto, de cualquier sucesión acotada en $H$ se puede extraer una subsucesión débilmente convergente.
\end{coro}

\section{Espacios uniformemente convexos}

\begin{definicion}
    Un espacio normado $X$ se dice \textbf{estrictamente convexo} si
    \begin{equation*}
        \forall x, y \in X \text{ con } x \neq y, \|x\| = \|y\| = 1 \implies \left\| \frac{x+y}{2} \right\| < 1
    \end{equation*}
    \textit{\small (el punto medio de 2 vectores unitarios cae dentro de la bola)}
\end{definicion}

\begin{ejemplo}
    $\ell_p, p > 1$ es estrictamente convexo, veámoslo
    \begin{equation*}
        \left\| \frac{1}{2}(x+y) \right\|_{\ell_p} < \frac{1}{2}(\|x\| + \|y\|) \quad \forall x, y \in \ell_p \text{ con } x \neq \alpha y 
    \end{equation*}
    \textit{\small (desigualdad estricta)}
\end{ejemplo}

\observacion{
    Veámos algunos ejemplos
    \begin{itemize}
        \item $(\mathbb{R}^n, d_2)$ es estrictamente convexo.
        \item $(\mathbb{R}^n, d_1)$ no es estrictamente convexo: 
        $$\left\| \frac{1}{2}(x+y) \right\| = \frac{1}{2}\|x\| + \frac{1}{2}\|y\| = 1 \text{ con } \|x\| = \|y\| = 1$$
    \end{itemize}
}

\begin{ejercicio}
    $\ell_1, \ell_\infty, C^0([0,1]), L^1(A)$ no son estrictamente convexos.
\end{ejercicio}

\begin{definicion}
    Un espacio normado $X$ se dice \tbf{uniformemente convexo} si 
    \begin{equation*}
        \forall \varepsilon > 0, \exists \delta > 0 \text{ tal que } \|x\|, \|y\| \leq 1, \|x - y\| \ge \varepsilon \implies \left\| \frac{x+y}{2} \right\| \leq 1 - \delta
    \end{equation*}
    \textit{\small (el punto medio de 2 vectores unitarios que están al menos $\varepsilon$ de distancia cae dentro de la bola cerrada de radio $1 - \delta$)}
\end{definicion}

\begin{lema}
    Sea $X$ un espacio de Banach, entonces
    \begin{equation*}
        X \text{ es reflexivo} \iff \forall F \in X^* \quad \exists x_0 \in X : \|x_0\| = 1 \text{ y } \|F\| = |F(x_0)|
    \end{equation*}
\begin{proof}
    \begin{description}
        \item[$\Longrightarrow$)] Ok.
        \item[$\Longleftarrow$)] Difícil, demostrado en 1970.
    \end{description}
\end{proof}
\end{lema}

\begin{teo}[Teorema de Milman-Pettis]
    Sea $X$ un espacio de Banach uniformemente convexo, entonces $X$ es reflexivo.
\begin{proof}   
    Sea $F \in X^* : \|F\| = 1$ sabemos que $\exists \{x_m\} \subseteq X$ con $\|x_m\| = 1$, simplemente se trata de construir la sucesión con los $x_m$ que hacen que $F(x_m)$ se acerque a $\|F\|$. De hecho, como $\|F\| = 1$, entonces
    \begin{equation*}
        \|F\|=\sup_{\|x\|=1} |F(x)| \ge |F(x_m)| \quad \forall m \implies \lim_{m \to +\infty} |F(x_m)-1| \leq  \lim_{m \to +\infty} \left|\sup_{\|x\|=1} |F(x)| -1\right|=\lim_{m \to +\infty} 0=0
    \end{equation*}
    por lo que $F(x_m) \longrightarrow 1 = \|F\|$. Demostremos que $\{x_m\}$ es de Cauchy, usando la reducción al absurdo, es decir, supongamos que no es de Cauchy, entonces
    \begin{equation*}
        \exists \bar{\varepsilon} > 0 : \forall k > 0, \exists m_k, n_k : \|x_{m_k} - x_{n_k}\| > \bar{\varepsilon}
    \end{equation*}
    Por la convexidad uniforme
    \begin{equation*}
        \exists \delta > 0 : \left\| \frac{1}{2}(x_{m_k} + x_{n_k}) \right\| \le 1 - \delta < 1 \implies |F(\frac{1}{2}(x_{m_k} + x_{n_k}))| \le \|F\| \cdot \left\| \frac{1}{2}(x_{m_k} + x_{n_k}) \right\| \le 1 - \delta < 1
    \end{equation*}
    lo cual es un absurdo, pues $\|F\|=1$, de hecho tenemos que 
    \begin{equation*}
        F(x_m) \longrightarrow \|F\| = 1 \implies \frac{1}{2}(F(x_{m_k}) + F(x_{n_k})) \longrightarrow 1
    \end{equation*}
    Demostrado queda entonces que $\{x_m\}$ es de Cauchy, y como $X$ es de Banach (completo), entonces $\{x_m\}$ converge a $\tilde{x} \in X$, con  $\|\tilde{x}\| = \lim\limits_{m \to \infty} \|x_m\| = 1$ y $F(x_m) \longrightarrow F(\tilde{x}) = 1$. Por lo tanto, $F$ alcanza su norma $\|F\| = 1$. \\ 

    \noindent
    Es claro que aqui no hemos usado un $F\in X^*$ cualquiera, sino uno con $\|F\|=1$, pero si $F$ es un funcional cualquiera, tenemos estos dos casos extras
    \begin{itemize}
        \item $F = 0$, entonces $F$ alcanza su norma en cualquier $x_0 \in X$ con $\|x_0\|=1$.
        \item $F \neq 0$ y $\|F\| \neq 0$, entonces  definimos 
        $$G = \frac{F}{\|F\|}$$ 
        un funcional con norma 1. Como nuestra demostración de Milman-Pettis funciona para cualquier funcional de norma 1, sabemos que existe un $\tilde{x} \in X$ con $\|\tilde{x}\|=1$ tal que $|G(\tilde{x})| = 1$. \\
        \noindent
        Ahora, simplemente ``deshacemos'' la división. Si $|G(\tilde{x})| = 1$, entonces:$$\left| \frac{F(\tilde{x})}{\|F\|} \right| = 1 \implies |F(\tilde{x})| = \|F\|$$
    \end{itemize}
    Y claramente aplicando el lema anterior, $X$ es reflexivo.
\end{proof}
\end{teo}


\section{Topologías débiles}
\noindent
Sea $X$ un conjunto y sea $\{Y_i\}$ una familia de espacios topológicos tales que existen funciones $\varphi_i: X \to Y_i$ para todo $i$. \\

\noindent
Queremos dotar a $X$ de una topología $\mathcal{T}$ que haga continuas a las $\{\varphi_i\}$ y que sea la menos fina (más pequeña) con esta propiedad. Es decir

\begin{equation*}
    \forall W_i\subseteq Y_i \text{ abierto}, \varphi_i^{-1}(W_i) \text{ es abierto en } (X, \mathcal{T}) \text{ para todo } i
\end{equation*}
\noindent
Sea $\mathcal{F}$ una familia de subconjuntos de $X$ tal que la intersección finita $\bigcap\limits_{\text{finitas}} \varphi_i^{-1}(W_i)$ está contenida en $X$ y contiene la unión arbitraria $\bigcup \varphi_i^{-1}(W_i)$. \\ \\
\noindent
Después de haber realizado estas dos operaciones obtenemos que $\mathcal{F}$ es estable por intersecciones finitas y uniones, y contiene a $\varphi_i^{-1}(W_i)$ para todo $W_i$ abierto en $Y_i$ y para todo $i$; por construcción, es la topología más pequeña con estas propiedades. \\
\noindent
Definimos entonces $\mathcal{T} = \mathcal{F}$.


\begin{prop}
    Sea $x_0 \in X$ y sea $\{\bigcap\limits_{\text{finitas}} \varphi_i^{-1}(V_i)\}$ una base de entornos de $x_0$, donde cada $V_i$ es un entorno de $\varphi_i(x_0)$ en $Y_i$ para todo $i$. Entonces
    \begin{equation*}
        x_n \to x \text{ en } (X, \mathcal{T}) \iff \varphi_i(x_n) \to \varphi_i(x) \text{ en } Y_i \quad \forall i
    \end{equation*}
\begin{proof}
    \begin{description}
        \item[$\Longrightarrow)$] Si $x_n \to x$, por la definición de $\mathcal{T}$, $\varphi_i$ es continua y por lo tanto $\varphi_i(x_n) \to \varphi_i(x)$.  
        \item[$\Longleftarrow)$] Sea $\mathcal{I}(x)$ un entorno de $x \in (X, \mathcal{T})$, entonces $\mathcal{I}(x) = \bigcap\limits_{\text{finitas}} \varphi_i^{-1}(V_i)$ con $V_i$ entorno de $\varphi_i(x)$. Por hipótesis $\varphi_i(x_n) \to \varphi_i(x)$, por lo que $\varphi_i(x_n) \in V_i$ para todo $n > \nu_i$ para algún $\nu_i > 0$. Sea $\nu = \max_i \{\nu_i\} < +\infty$, entonces para todo $n > \nu$ se tiene que $\varphi_i(x_n) \in V_i$ para todo $i$ y, por lo tanto, $x_n \in \varphi_i^{-1}(V_i)$ para todo $i$. \\ \\
        \noindent
        Esto implica que $x_n \in \bigcap_i \varphi_i^{-1}(V_i) = \mathcal{I}(x) \implies x_n \to x$. 
    \end{description}
\end{proof}
\end{prop}

\begin{definicion}
    Sea $X$ un espacio normado y $\mathcal{T}$ la topología menos fina que hace continuas las aplicaciones $F \in X^*$. Definimos $\mathcal{T}$ como la \textbf{topología débil}. \textit{\small (Es el caso precedente con $Y_i = \mathbb{R} \ \forall i$)}
\end{definicion}

\observacion{
    $(X, \|\cdot\|)$ hace continuas las aplicaciones de $X^*$ y, por lo tanto, también las de $(X, \mathcal{T})$.
    \[ x_m \xrightarrow{\mathcal{T}} x \text{ en } (X, \mathcal{T}) \iff F(x_m) \longrightarrow F(x) \text{ en } \mathbb{R} \ \forall F \in X^* \iff x_m \rightharpoonup x \]

}

\observacion{
    Sea $x_0 \in X$, una base de entornos se obtiene como
    \[ V_\varepsilon(x_0) = \{ x \in X \mid |F(x) - F(x_0)| < \varepsilon \ \forall F \text{ en un subconjunto finito de } X^* \} \]
}

\begin{prop}
    $(X, \mathcal{T})$ satisface el axioma de separación $T_2$ (Hausdorff).
\begin{proof}
    Hahn-Banach geométrico sobre los funcionales.
\end{proof}
\end{prop}

\observacion{
    En dimensión finita, la topología fuerte y la topología débil coinciden.
}

\begin{observacion}
    Cuando hacemos alusión a la topología fuerte en un espacio normado $(X,\|\cdot\|)$, nos referimos a la topología inducida por la norma $\|\cdot\|$ y, por lo tanto, a la convergencia fuerte. Por otro lado, cuando hablamos de la topología débil en $(X, \mathcal{T})$, nos referimos a la topología menos fina que hace continuas a las aplicaciones de $X^*$ y, por lo tanto, a la convergencia débil.
\end{observacion}

\observacion{
    Sea $(X, \|\cdot\|)$ un espacio normado de dimensión infinita, entonces $(X, \mathcal{T})$ débil es siempre estrictamente menos fina (más pequeña) que $(X, \|\cdot\|)$ con topología fuerte. En particular
    \begin{equation*}
        B = \{ x \in X \mid \|x\| < 1 \}
    \end{equation*}
    nunca es abierto en $(X, \mathcal{T})$, pues se puede ver que cada entorno de $x_0 \in (X, \mathcal{T})$ contiene una recta vectorial que pasa por $x_0$, por lo tanto, $B$ nunca es abierto. 
}

\observacion{
    En dimensión infinita, $(X, \|\cdot\|)$ nunca es metrizable. Existen espacios en los que la convergencia por sucesiones 
    $$x_m \to x \iff x_m \rightharpoonup x$$
    pero no es metrizable. 
}

\begin{ejercicio}
    En $\ell_1$, la convergencia fuerte y la convergencia débil coinciden, es decir
    \begin{equation*}
            x_m \to x \text{ en } \ell_1 \iff x_m \rightharpoonup x \text{ en } \ell_1
    \end{equation*}
\end{ejercicio}

\noindent
Intuitivamente, la topología débil tiene ``menos abiertos'' y, por lo tanto, ``más compactos'', entonces podemos minimizar funcionales sobre los compactos.